%% kapa: exam=1 round=33 year=2022 period=1 form=A subject=civil_law subject_ko=민법 seq=18 number=18
%% source: Q-Net 감정평가사 기출문제 (cst00309) artlSeq=5209997 file=raw/1008/5209997/2022년도 제33회 감정평가사 1차시험 1교시(민법,경제학,부동산학).pdf
%% conversion: kordoc 4.24.51 -> tools/gen1.py (auto); review=proofread-sonnet; flags=none
\begin{question}{18}
법률행위의 조건에 관한 설명으로 옳지 않은 것은?
\begin{choices}
  \item 조건의 성취로 인하여 이익을 받을 당사자가 신의성실에 반하여 조건을 성취시킨 때에는 상대방은 그 조건이 성취하지 아니한 것으로 주장할 수 있다.
  \item 법률행위 당시 이미 성취된 조건을 해제조건으로 하는 법률행위는 조건 없는 법률행위이다.
  \item 정지조건이 있는 법률행위는 특별한 사정이 없는 한 조건이 성취한 때로부터 그 효력이 생긴다.
  \item 조건 있는 법률행위의 당사자는 조건의 성부가 미정한 동안에 조건의 성취로 인하여 생길 상대방의 이익을 해하지 못한다.
  \item 조건의 성취가 미정인 권리도 일반규정에 의하여 담보로 할 수 있다.
\end{choices}
\end{question}
\answer{2}
